%%%%%%%%%%%%%%%%%%%%%%%%%%%%%%%%%%%%%%%%%%%%%%%%%%%%%%%%%%%%%%%%%%%%%%%%%%%%%%%%
%\section*{Appendix: Target-Mass Correction for Quasi-Distributions}
%%%%%%%%%%%%%%%%%%%%%%%%%%%%%%%%%%%%%%%%%%%%%%%%%%%%%%%%%%%%%%%%%%%%%%%%%%%%%%%%

In this Appendix, we derive the target-mass corrections to the unpolarized and helicity parton distributions. 

%%%%%%%%%%%%%%%%%%%%%%%%%%%%%%%%%%%%%%%%%%%%%%%%%%%%%%%%%%%%%%%%%%%%%%%%%%%%%%%%
\subsection{Unpolarized distribution}

For the unpolarized parton distribution, the series sum in Eq.~\ref{Kn} can be explicitly performed, and the result for an even $n$ ($=2k$) is
%
\begin{equation}\label{evensum}
\sum_{j=0}^k  C_{n-j}^j c^j =\frac{1}{\sqrt{1+4c}}\Big[\left(\frac{f_-}{2}\right)^{2k+1}+\left(\frac{f_+}{2}\right)^{2k+1}\Big],
\end{equation}
%
and for an odd $n$ ($=2k+1$) is
%
\begin{equation}\label{oddsum}
\sum_{j=0}^k C_{n-j}^j c^j =\frac{1}{\sqrt{1+4c}}\Big[-\left(\frac{f_-}{2}\right)^{2k+2}+\left(\frac{f_+}{2}\right)^{2k+2}\Big].
\end{equation}
%
where $f_{\pm}=\sqrt{1+4c}\pm 1$. 

With Eqs.~\ref{evensum} and \ref{oddsum}, we perform an inverse Mellin transform to the moment relation of Eq.~\ref{Kn}
%
\begin{equation}
\frac{1}{2\pi i}\int_{-i\infty}^{i\infty} dn \, s^{-n} \langle x^{n-1} \rangle .
\end{equation}
%
For $n=2k$, we obtain
%
\begin{equation}\label{pdfrel}
\tilde q(x)-\tilde q(-x)=\frac{1}{\sqrt{1+4c}}\sum_{i=\pm}\frac{f_i}{2}\Big[q\Big(\frac{2x}{f_i}\Big)-q\Big(\frac{-2x}{f_i}\Big)\Big],
\end{equation}
%
and for $n=2k+1$, we have
%
\begin{equation}\label{oddpdfrel}
\tilde q(x)+\tilde q(-x)=\frac{1}{\sqrt{1+4c}}\Big\{\frac{f_+}{2}\Big[q\Big(\frac{2x}{f_+}\Big)+q\Big(\frac{-2x}{f_+}\Big)\Big]-\frac{f_-}{2}\Big[q\Big(\frac{2x}{f_-}\Big)+q\Big(\frac{-2x}{f_-}\Big)\Big]\Big\},
\end{equation}
%
where we have used the following representation for the Dirac $\delta$-function~\cite{Steffens:2012jx}
%
\begin{equation}
\delta(\ln u)=\frac{1}{2\pi i}\int_{-i\infty}^{i\infty}dn\, u^n.
\end{equation}
%
From the above two equations, we have
%
\begin{equation}\label{qinq0}
\tilde q(x)=\frac{1}{\sqrt{1+4c}}\Big[\frac{f_+}{2}q\Big(\frac{2x}{f_+}\Big)-\frac{f_-}{2}q\Big(\frac{-2x}{f_-}\Big)\Big].
\end{equation}
%
Note that this is different from the result obtained in Ref.~\cite{Alexandrou:2015rja}. Their result does not conserve quark number. In our case it is easy to check the quark-number conservation
%
\begin{equation}
\int_{-\infty}^{\infty} dx\, \tilde q(x)=\frac{1}{\sqrt{1+4c}}\int dx \big[\frac{f_+^2}{4}-\frac{f_-^2}{4}\big]q(x)=\int_{-\infty}^{\infty} dx\, q(x),
\end{equation}
%
or $\langle x^0\rangle_{\tilde q}=\langle x^0\rangle_q$.

Eq.~\ref{qinq0} writes $\tilde q(x)$ in terms of $q(x)$, where the former is the quantity that can be directly computed on the lattice and the latter is the usual parton distribution. In practice, we would like to extract $q(x)$ from $\tilde q(x)$. To see how this can be done, let us rewrite Eq.~\ref{Kn} for an even $n=2k$ as
%
\begin{equation}\label{momentinverserel}
\langle x^{2k-1}\rangle_q=\langle x^{2k-1}\rangle_{\tilde q}\frac{\sqrt{1+4c}}{\left(\frac{f_-}{2}\right)^{2k+1}+\left(\frac{f_+}{2}\right)^{2k+1}}=\langle x^{2k-1}\rangle_{\tilde q}\frac{\sqrt{1+4c}}{\left(\frac{f_+}{2}\right)^{2k+1}}\sum_{n=0}^\infty (-1)^n\left(\frac{f_-}{f_+}\right)^{(2k+1)n}.
\end{equation}
%
The inverse Mellin transform then leads to
%
\begin{align}
q(x)-q(-x)=\sqrt{1+4c}\sum_{n=0}^\infty \frac{2(-f_-)^n}{f_+^{n+1}}\Big[\tilde q\Big(\frac{f_+^{n+1}x}{2f_-^n}\Big)-\tilde q\Big(\frac{-f_+^{n+1}x}{2f_-^n}\Big)\Big].
\end{align}
%
Similarly, we have
%
\begin{equation}
q(x)+q(-x)=\sqrt{1+4c}\sum_{n=0}^\infty \frac{2 f_-^n}{f_+^{n+1}}\Big[\tilde q\Big(\frac{f_+^{n+1}x}{2f_-^n}\Big)+\tilde q\Big(\frac{-f_+^{n+1}x}{2f_-^n}\Big)\Big].
\end{equation}
%
Therefore,
%
\begin{align}
q(x)&=\sqrt{1+4c}\sum_{n=0}^\infty \frac{f_-^n}{f_+^{n+1}}\Big[(1+(-1)^n)\tilde q\Big(\frac{f_+^{n+1}x}{2f_-^n}\Big)+(1-(-1)^n)\tilde q\Big(\frac{-f_+^{n+1}x}{2f_-^n}\Big)\Big]\non\\
&=\sqrt{1+4c}\sum_{n=0}^\infty \frac{(4c)^n}{f_+^{2n+1}}\Big[(1+(-1)^n)\tilde q\Big(\frac{f_+^{2n+1}x}{2(4c)^n}\Big)+(1-(-1)^n)\tilde q\Big(\frac{-f_+^{2n+1}x}{2(4c)^n}\Big)\Big],
\end{align}
%
where in the last line we have used $f_+ f_-=4c$. Since $f_+\gg f_-$ or $c$ and the quasi-distribution $\tilde q(x)$ vanishes asymptotically for large $x$, the above sum is dominated by the first term with $n=0$. In practical calculations, we can reach a reasonable accuracy by taking only the first few terms in the sum.

%%%%%%%%%%%%%%%%%%%%%%%%%%%%%%%%%%%%%%%%%%%%%%%%%%%%%%%%%%%%%%%%%%%%%%%%%%%%%%%%
\subsection {Helicity distribution}

Now let us look at the quark helicity distribution. In this case, Eq.~\ref{Knbar} gives the following relation between moments
%
\begin{equation}
\langle x^{n-1}\rangle_{\Delta q} \frac{1}{n}\sum_{j=0}^{\infty}c^j\frac{(n-j)!}{j!(n-2j-1)!}=\langle x^{n-1}\rangle_{\Delta \tilde q} .
\end{equation}
% 
The result of the series sum for an even $n=2k$ is
%
\begin{equation}
\sum_{j=0}^{k-1}c^j\frac{(n-j)!}{j!(n-2j-1)!}=\frac{1}{1+4c}\Big[\Big(2k+1 -\frac{1}{\sqrt{1+4c}}\Big)\left(\frac{f_+}{2}\right)^{2k+1}-\Big(2k+1 +\frac{1}{\sqrt{1+4c}}\Big)\left(\frac{f_-}{2}\right)^{2k+1}\Big],
\end{equation}
%
and for an odd $n=2k+1$ is
%
\begin{equation}
\sum_{j=0}^{k}c^j\frac{(n-j)!}{j!(n-2j-1)!}=\frac{1}{1+4c}\Big[\Big(2k+2 -\frac{1}{\sqrt{1+4c}}\Big)\left(\frac{f_+}{2}\right)^{2k+2}+\Big(2k+2 +\frac{1}{\sqrt{1+4c}}\Big)\left(\frac{f_-}{2}\right)^{2k+2}\Big].
\end{equation}
%
From the above equations, we have
%
\begin{align}
2k \Big[\langle x^{2k-1}\rangle_{\Delta \tilde q}-\frac{\langle x^{2k-1}\rangle_{\Delta q}}{1+4c}\Big(\left(\frac{f_+}{2}\right)^{2k+1}-\left(\frac{f_-}{2}\right)^{2k+1}\Big)\Big]&=\frac{2c \langle x^{2k-1}\rangle_{\Delta q}}{(1+4c)^\frac{3}{2}}\Big[\left(\frac{f_+}{2}\right)^{2k}-\left(\frac{f_-}{2}\right)^{2k}\Big], \non\\
(2k+1) \Big[\langle x^{2k}\rangle_{\Delta \tilde q}-\frac{\langle x^{2k}\rangle_{\Delta q}}{1+4c}\Big(\left(\frac{f_+}{2}\right)^{2k+2}+\left(\frac{f_-}{2}\right)^{2k+2}\Big)\Big]&=\frac{2c \langle x^{2k}\rangle_{\Delta q}}{(1+4c)^\frac{3}{2}}\Big[\left(\frac{f_+}{2}\right)^{2k+1}+\left(\frac{f_-}{2}\right)^{2k+1}\Big].
\end{align}
%
Their inverse Mellin transform then leads to
%
\begin{align}
&-\frac{\partial}{\partial s}\Big\{\frac{|s|}{s}\Big[\Delta\tilde q(s)-\Delta\tilde q(-s)-\frac{1}{1+4c}\left(\frac{f_+}{2}\Big(\Delta q\Big(\frac{2s}{f_+}\Big)-\Delta q\Big(\frac{-2s}{f_+}\Big)\Big)-\frac{f_-}{2}\Big(\Delta q\Big(\frac{2s}{f_-}\Big)-\Delta q\Big(\frac{-2s}{f_-}\Big)\Big)\right)\Big]\Big\}\non\\
&=\frac{2c}{(1+4c)^{\frac{3}{2}}|s|}\Big\{\Big[\Delta q\Big(\frac{2s}{f_+}\Big)-\Delta q\Big(\frac{-2s}{f_+}\Big)\Big]-\Big[\Delta q\Big(\frac{2s}{f_-}\Big)-\Delta q\Big(\frac{-2s}{f_-}\Big)\Big]\Big\},\non\\
&-\frac{\partial}{\partial s}\Big\{\frac{|s|}{s}\Big[\Delta \tilde q(s)+\Delta \tilde q(-s)-\frac{1}{1+4c}\left(\frac{f_+}{2}\Big(\Delta q\Big(\frac{2s}{f_+}\Big)+\Delta q\Big(\frac{-2s}{f_+}\Big)\Big)+\frac{f_-}{2}\Big(\Delta q\Big(\frac{2s}{f_-}\Big)+\Delta q\Big(\frac{-2s}{f_-}\Big)\Big)\Big)\right)\Big]\Big\}\non\\
&=\frac{2c}{(1+4c)^{\frac{3}{2}}|s|}\Big\{\Big[\Delta q\Big(\frac{2s}{f_+}\Big)+\Delta q\Big(\frac{-2s}{f_+}\Big)\Big]+\Big[\Delta q\Big(\frac{2s}{f_-}\Big)+\Delta q\Big(\frac{-2s}{f_-}\Big)\Big]\Big\}.
\end{align}
%
From these two equations, we obtain
%
\begin{equation}
-\frac{\partial}{\partial x}\Big[\Delta \tilde q(x)-\frac{1}{1+4c}\Big(\frac{f_+}{2}\Delta q\Big(\frac{2x}{f_+}\Big)+\frac{f_-}{2}\Delta q\Big(\frac{-2x}{f_-}\Big)\Big)\Big]=\frac{2c}{(1+4c)^{\frac{3}{2}}x}\Big[\Delta q\Big(\frac{2x}{f_+}\Big)+\Delta q\Big(\frac{-2x}{f_-}\Big)\Big].
\end{equation}
%
Its solution is given by
%
\begin{equation}
\Delta \tilde q(x)=\frac{1}{1+4c}\Big(\frac{f_+}{2}\Delta q\Big(\frac{2x}{f_+}\Big)+\frac{f_-}{2}\Delta q\Big(\frac{-2x}{f_-}\Big)\Big)-\int_{\pm\infty}^x \frac{dy}{y}\frac{2c}{(1+4c)^{\frac{3}{2}}}\Big[\Delta q\Big(\frac{2y}{f_+}\Big)+\Delta q\Big(\frac{-2y}{f_-}\Big)\Big],
\end{equation}
%
where we have used the fact that the quasi-distribution vanishes sufficiently fast for $x\to\pm\infty$. It is irrelevant to choose $\infty$ or $-\infty$ in the lower limit of the second integral. To facilitate numerical implementation, we choose it as $\infty$ ($-\infty$) for $x>0$ ($x<0$).

Since $f_+\gg f_-$ or $c$, we can recursively solve the above equation for $\Delta q(x)$ by making a change of variable $x\to f_+ x/2, y\to f_+ y/2$. 
For simplicity, let us denote
%
\begin{equation}
a=1+4c, \hspace{2em} b=\frac{f_+}{\sqrt{1+4c}}, \hspace{2em} r=\frac{f_-}{f_+}, \hspace{2em} \Delta(y)=\Delta \tilde q\left(\frac{f_+}{2}y\right)+\Delta \tilde q\left(-\frac{f_+}{2}\frac{y}{r}\right).
\end{equation}
%
We then have
%
\begin{align}
\Delta q(x<0)&=\frac{2a}{f_{+}}\Bigg\{\Delta \tilde{q}\left( \frac{f_{+}}{2}%
x\right) -r\left[\Delta \tilde{q}\left( -\frac{f_{+}}{2}\frac{x}{r}\right)
-\int_{-\infty }^{x}\frac{dy}{y}\,b\,\Delta (y)\right]\non \\
& +r^{2}\left[\Delta \tilde{q}\left( \frac{f_{+}}{2}\frac{x}{r^{2}}\right)
-\int_{-\infty }^{x}\frac{dy}{y}\,b\,\Delta \Big(-\frac{y}{r}\Big)%
+\int_{-\infty }^{x}\frac{dy}{y}\int_{-\infty }^{y}\frac{dz}{z}\,b^{2}\Delta
(z)-\int_{-\infty }^{-\frac{x}{r}}\frac{dy}{y}\,b\,\Delta (y)\right.\non \\
& \left.+\int_{-\infty }^{x}\frac{dy}{y}\int_{-\infty }^{-\frac{y}{r}}\frac{dz}{z}%
\,b^{2}\Delta (z)\right]\\
& -r^3 \left[ \Delta \tilde{q}\left( -\frac{%
f_{+}}{2r^{3}}x\right) -b^{3}\left( \int_{-\infty }^{z}\frac{dk}{k}%
\int_{-\infty }^{y}\frac{dz}{z}\int_{-\infty }^{x}\frac{dy}{y}\Delta
(k)+\int_{-\infty }^{z}\frac{dk}{k}\int_{-\infty }^{-\frac{y}{r}}\frac{dz}{z}%
\int_{-\infty }^{x}\frac{dy}{y}\Delta (k)\right. \right.  \\
& +\left. \int_{-\infty }^{-\frac{z}{r}}\frac{dk}{k}\int_{-\infty }^{y}\frac{%
dz}{z}\int_{-\infty }^{x}\frac{dy}{y}\Delta (k)+\int_{-\infty }^{-\frac{z}{r}%
}\frac{dk}{k}\int_{-\infty }^{-\frac{y}{r}}\frac{dz}{z}\int_{-\infty }^{x}%
\frac{dy}{y}\Delta (k)\right)  \\
& +b^{2}\left( \int_{-\infty }^{\frac{y}{r^{2}}}\frac{dk}{k}\int_{-\infty
}^{x}\frac{dy}{y}\Delta (k)+\int_{-\infty }^{-\frac{y}{r}}\frac{dk}{k}%
\int_{-\infty }^{x}\frac{dy}{y}\Delta (k)+\int_{-\infty }^{y}\frac{dk}{k}%
\int_{-\infty }^{-\frac{x}{r}}\frac{dy}{y}\Delta (k)\right.  \\
& \left. +\int_{-\infty }^{-\frac{y}{r}}\frac{dk}{k}\int_{-\infty }^{-\frac{x%
}{r}}\frac{dy}{y}\Delta (k)+\int_{-\infty }^{\frac{y}{r}}\frac{dk}{k}%
\int_{-\infty }^{x}\frac{dy}{y}\Delta (-k)+\int_{-\infty }^{-\frac{y}{r^{2}}}%
\frac{dk}{k}\int_{-\infty }^{x}\frac{dy}{y}\Delta (-k)\right)  \\
& \left.\left. -b\left( 2\int_{-\infty }^{\frac{x}{r^{2}}}\frac{dk}{k}\Delta
(k)+\int_{-\infty }^{-\frac{x}{r^{2}}}\frac{dk}{k}\Delta (-k)\right) \right] \right\}
 +\mathcal{O}(r^{4}),\non\\
\Delta q(x>0)&=\Delta q(x<0)[-\infty\to\infty].
\end{align}
%
We have derived the mass correction up to $\mathcal O(r^3)$. Although in the present work we did not implement $\mathcal O(r^3)$ correction due to its computational complexity, the above result will be useful for future improvements with more computational resources.
